\chapter{BZN\({}^{\mathrm{Ar}}\): Loop-Space Convention and Wild-Prime Summands}
\label{v4-arith:bzn-rcompletion}

The F3.BZN\({}^{\mathrm{Ar}}\) reduction of
Chapter~\ref{v4-arith:arith-bzn} rests on two unresolved conditions.
The first: the arithmetic free loop space
\(L^{\mathrm{Ar}}\mathcal{X}=\mathrm{Map}(B\widehat{\mathbb{Z}},\mathcal{X})\)
admits three candidate definitions on formal Noetherian derived
stacks, differing a priori in the formal-Noetherian setting; no
canonical choice was specified. The second: the Iwahori Hecke
category \(\mathrm{Hk}^{\mathrm{Iw}}_{GL_{2},p}\) is not semisimple at
the wild primes \(p=2,3\) of \(GL_{2}/\mathbb{Q}\), so the Hochschild
trace picks up \(\mathrm{Ext}^{1}\)-contributions invisible to the
Ben-Zvi--Nadler semisimple template.
The derived internal-Hom convention is the convention used below.
Agreement with the cofiltered-limit and ind-construction conventions is
a formal-convergence hypothesis on the target, not a theorem of the
definition. The wild-prime condition is modeled by an explicit
calculation whose conclusion depends on unverified
coefficient-convention compatibility: the expected
\(\mathrm{Ext}^{1}\)-summand at \(p=2\) and \(p=3\) is computed via
Etingof--Ostrik 2004 \emph{Finite Tensor Categories},
Bushnell--Henniart 2006 local Langlands at the wild primes, and
Jantzen 2003 modular-representation Ext-machinery; the coefficient
conventions must be verified against the arithmetic mixed-characteristic
setting before the calculation has independent source support.
The additive wild summand required by the proposed
F3.BZN\({}^{\mathrm{Ar}}\) comparison is recorded. The global
F3.BZN\({}^{\mathrm{Ar}}\) statement of Chapter~\ref{v4-arith:arith-bzn}
remains conditional until the loop-space convergence, full
Weil/\((\varphi,\Gamma)\) target, wild-prime Ext, and \(p\geq 5\)
semisimplicity inputs are verified in the arithmetic coefficient
conventions.

\section{The Two Loop-Convention and Wild-Prime Conditions}
\label{v4-arith:bzn-rcompletion:context}

\subsection{Origin of the Two Conditions}

The structural comparison of Chapter~\ref{v4-arith:comparison} at
\S\ref{v4-arith:sw:wave2} contains four conditions in the F3.BZN\({}^{\mathrm{Ar}}\)
reduction. Three of these (the formal-Noetherian
extension of Gaitsgory--Rozenblyum; the cotangent-complex check for
the Emerton--Gee stack; and the R1 relabeling) were absorbed into
the R1 split \(\mathrm{R1.fin\text{-}type}+\mathrm{R1.formal}\)
stated in Chapter~\ref{v4-arith:comparison}. The remaining two,
not absorbed into the existing reduction, are:
\begin{itemize}
\item \textbf{R3.} Ambiguity in the loop-space convention.
Three candidate definitions of
\(L^{\mathrm{Ar}}\mathcal{X}=\mathrm{Map}(B\widehat{\mathbb{Z}},\mathcal{X})\)
are in play:
\begin{itemize}
\item[\textbf{(a)}] \emph{Cofiltered-limit convention.}
\(L^{\mathrm{Ar}}_{\mathrm{cofilt}}\mathcal{X}:=\varprojlim_{n}\mathrm{Map}(B(\mathbb{Z}/n),\mathcal{X})\),
the limit over finite quotients of \(\widehat{\mathbb{Z}}\).
\item[\textbf{(b)}] \emph{Internal-Hom convention.}
\(L^{\mathrm{Ar}}_{\mathrm{intHom}}\mathcal{X}:=\underline{\mathrm{Map}}_{\mathrm{dSt}^{\wedge}_{\mathrm{Noeth}}}(B\widehat{\mathbb{Z}},\mathcal{X})\),
the internal mapping stack in the \(\infty\)-category of formal
Noetherian derived stacks.
\item[\textbf{(c)}] \emph{Ind-construction convention.}
\(L^{\mathrm{Ar}}_{\mathrm{ind}}\mathcal{X}:=\varinjlim_{n}\mathrm{Ind}\bigl(\mathrm{Map}(B(\mathbb{Z}/n),\mathcal{X}^{\mathrm{perf}})\bigr)\),
the ind-completion over finite quotients, applied to the perfect
locus and re-ind-extended.
\end{itemize}
These agree for \(\mathcal{X}\) finite-type smooth, but the formal
Noetherian setting is sensitive to the order of limits and
ind-completions.

\item \textbf{R4.} Non-semisimplicity at the wild primes.
For \(GL_{2}/\mathbb{Q}\) the non-banal primes are \(p=2,3\): at these
primes the Iwahori Hecke category
\(\mathrm{Hk}^{\mathrm{Iw}}_{GL_{2},p}\) is not semisimple. The
classical Ben-Zvi--Nadler theorem, and its R1/R2/R3 reduction,
presupposes semisimplicity at every factor; the Hochschild trace
over a non-semisimple factor involves higher \(\mathrm{Ext}^{i}\)
terms, which the ``\(\mathcal{O}(L\mathcal{X})^{\mathrm{flat}}\)''
model does not see.
\end{itemize}

\subsection{Logical Structure}

R3 is a convention plus a convergence hypothesis: the internal Hom is
the chosen definition, and comparison with the other two definitions
requires compactness and formal-colimit compatibility. R4 is a
conditional wild-Ext model: the \(p=2,3\) summands and the \(p\geq5\)
vanishing are hypotheses until the coefficient conventions and
cyclic-trace pairings are verified.

\subsection{Notation}

Throughout this chapter:
\begin{itemize}
\item \(\widehat{\mathbb{Z}}=\prod_{p}\mathbb{Z}_{p}\), the profinite
completion of \(\mathbb{Z}\).
\item \(\mathrm{dSt}^{\wedge}_{\mathrm{Noeth}}\): the
\(\infty\)-category of formal Noetherian derived (pro-\'etale)
stacks in the sense of Gaitsgory--Rozenblyum's AMS DAG monographs
Vol.~II Part~III. An object is a pro-object \(\{\mathcal{X}_{n}\}\)
of derived Noetherian stacks with \(\mathcal{X}_{n}\) classically a
Noetherian topological ring with derived enhancement.
\item For \(\mathcal{X}\in\mathrm{dSt}^{\wedge}_{\mathrm{Noeth}}\),
its formal completion \(\mathcal{X}^{\wedge}\) is the image of
\(\mathcal{X}\) under the reflective embedding into the
formal-Noetherian \(\infty\)-category.
\item \(\mathrm{Hk}^{\mathrm{Iw}}_{GL_{2},p}\): the Iwahori Hecke
category at the prime \(p\) (Bezrukavnikov arXiv:1209.0403).
\item \(\mathrm{char}(k)=\ell\) denotes the characteristic of the
coefficient field \(k\); we typically take \(k=\overline{\mathbb{Q}}_{\ell}\)
with \(\ell\neq p\) but also address \(k=\overline{\mathbb{F}}_{\ell}\).
\item \(\mathrm{Ext}^{1}_{\mathcal{A}}\): \(\mathrm{Ext}^{1}\) in the
abelian category \(\mathcal{A}\); subscripts indicate the ambient
category.
\end{itemize}

\section{Structure of the Argument}
\label{v4-arith:bzn-rcompletion:table}

The argument for R3 has three parts; the argument for R4 has three
further parts. For R3: the three candidate definitions are compared
with explicit fibre-product formulas
(\Cref{v4-arith:bzn-rcompletion:prop:R3-three-computations}); the fibre-product formulas
are conditionally identified on formal Noetherian derived algebraic
stacks via Gaitsgory--Rozenblyum and Postnikov-tower convergence
(\Cref{v4-arith:bzn-rcompletion:thm:R3-three-agree}); the internal-Hom convention is
adopted as canonical and confirmed consistent with Chapter~\ref{v4-arith:arith-bzn}
(\Cref{v4-arith:bzn-rcompletion:def:canonical-L-Ar},
\Cref{v4-arith:bzn-rcompletion:prop:R3-matches-ch28}).
For R4: the \(\mathrm{Ext}^{1}\)-contribution to the Hochschild trace at \(p=2\) is
isolated via Bushnell--Henniart 2006 and Jantzen 2003
(\Cref{v4-arith:bzn-rcompletion:prop:R4-Ext1-p2}); the computation at \(p=3\) is
isolated via Carter 1989
(\Cref{v4-arith:bzn-rcompletion:prop:R4-Ext1-p3}); the wild summand vanishes at \(p\geq 5\)
only under the large-prime semisimplicity datum
(\Cref{v4-arith:bzn-rcompletion:thm:R4-correction-vanishes-p-large}).
Each step carries its source and convention hypotheses inline.

\section{R3: Arithmetic Loop-Space Convention}
\label{v4-arith:bzn-rcompletion:R3}

\subsection{The Three Candidate Definitions}

Throughout this section \(\mathcal{X}\in\mathrm{dSt}^{\wedge}_{\mathrm{Noeth}}\)
is a formal Noetherian derived algebraic stack over \(\mathbb{Z}\)
(or over \(\mathrm{Spf}\,\mathbb{Z}_{p}\) when working locally at a
prime). We demand:
\begin{itemize}
\item \(\mathcal{X}\) is locally presented as the formal completion
of a derived Noetherian affine scheme
\(\mathcal{X}_{0}=\mathrm{Spec}\,A_{0}^{\bullet}\) along an ideal
\(I\subset\pi_{0}(A_{0}^{\bullet})\);
\item the cotangent complex \(\mathbb{L}_{\mathcal{X}}\) has finite
Tor-amplitude \([a,b]\) over the structure sheaf;
\item the Postnikov tower \(\pi_{\leq n}\mathcal{X}\to\mathcal{X}\)
stabilises in each tor-layer (equivalently, \(\mathcal{X}\) is
\(t\)-bounded in the sense of Gaitsgory--Rozenblyum Vol.~I~\S2).
\end{itemize}
For the arithmetic target these conditions must be checked
for the Selmer-derived stack
\(\mathcal{X}^{\mathrm{Sel}}_{\bar\rho,\det,\mathcal{L}}\); the local
Emerton--Gee theorem supplies only the local formal inputs.

\begin{definition}[the three candidate arithmetic free loop spaces]
\label{v4-arith:bzn-rcompletion:def:three-candidates}
For \(\mathcal{X}\in\mathrm{dSt}^{\wedge}_{\mathrm{Noeth}}\), define
\begin{align}
L^{\mathrm{Ar}}_{\mathrm{cofilt}}\mathcal{X}
&\;:=\;\varprojlim_{n\geq 1}\mathrm{Map}\bigl(B(\mathbb{Z}/n\mathbb{Z}),\mathcal{X}\bigr),\\
L^{\mathrm{Ar}}_{\mathrm{intHom}}\mathcal{X}
&\;:=\;\underline{\mathrm{Map}}_{\mathrm{dSt}^{\wedge}_{\mathrm{Noeth}}}\bigl(B\widehat{\mathbb{Z}},\mathcal{X}\bigr),\\
L^{\mathrm{Ar}}_{\mathrm{ind}}\mathcal{X}
&\;:=\;\varinjlim_{n\geq 1}\mathrm{Ind}\bigl(\mathrm{Map}(B(\mathbb{Z}/n\mathbb{Z}),\mathcal{X}^{\mathrm{perf}})\bigr),
\end{align}
where:
\begin{itemize}
\item in (a) the limit is along the cofiltered system
\(\{\mathbb{Z}/n\}_{n\in\mathbb{N}}\) under divisibility;
\item in (b) the internal mapping stack is computed in the
\(\infty\)-category \(\mathrm{dSt}^{\wedge}_{\mathrm{Noeth}}\) of
formal Noetherian derived pro-\'etale stacks;
\item in (c) \(\mathcal{X}^{\mathrm{perf}}\) denotes the perfect
locus (dualisable objects in the formal quasi-coherent
derived category), and \(\mathrm{Ind}\) is the ind-completion in
stable presentable \(\infty\)-categories.
\end{itemize}
\end{definition}

\subsection{Explicit Computation of the Three Candidate Definitions}

\begin{proposition}[three candidate loop-space formulas]
\label{v4-arith:bzn-rcompletion:prop:R3-three-computations}
\ClaimStatusConjectured. Let \(\mathcal{X}\in\mathrm{dSt}^{\wedge}_{\mathrm{Noeth}}\)
be formal Noetherian derived algebraic with cotangent complex
\(\mathbb{L}_{\mathcal{X}}\) of finite Tor-amplitude. Then:
\begin{enumerate}
\item \emph{Cofiltered-limit formula.}
\begin{equation}
L^{\mathrm{Ar}}_{\mathrm{cofilt}}\mathcal{X}
\;=\;\varprojlim_{n}\Bigl(\mathcal{X}\times_{\mathcal{X}\times\mathcal{X}}\mathcal{X}_{(n)}\Bigr),
\end{equation}
where \(\mathcal{X}_{(n)}\subset\mathcal{X}\times\mathcal{X}\) is the
\(n\)-th formal neighbourhood of the diagonal.
\item \emph{Internal-Hom formula.}
\begin{equation}
L^{\mathrm{Ar}}_{\mathrm{intHom}}\mathcal{X}
\;=\;\mathcal{X}\times_{\mathcal{X}\times\mathcal{X}}^{\mathrm{d},\wedge}\mathcal{X},
\end{equation}
the derived formal self-intersection of \(\mathcal{X}\) along its
diagonal, computed in \(\mathrm{dSt}^{\wedge}_{\mathrm{Noeth}}\)
with the derived tensor product taken over the formal diagonal.
\item \emph{Ind-construction formula.}
\begin{equation}
L^{\mathrm{Ar}}_{\mathrm{ind}}\mathcal{X}
\;=\;\mathrm{Ind}\Bigl(\varinjlim_{n}\mathcal{X}^{\mathrm{perf}}\times_{\mathcal{X}^{\mathrm{perf}}\times\mathcal{X}^{\mathrm{perf}}}\mathcal{X}^{\mathrm{perf}}_{(n)}\Bigr).
\end{equation}
\end{enumerate}
\end{proposition}

\begin{proof}[Evidence]
(1) The mapping stack from \(B(\mathbb{Z}/n)\) to \(\mathcal{X}\) is
the fixed-point locus of the order-\(n\) automorphism action, which
is the derived fibre product along the \(n\)-th iterate of the
diagonal. In characteristic zero (for the rational structure at
each prime) this reduces to the formal neighbourhood
\(\mathcal{X}_{(n)}\). Taking the limit over \(n\) under divisibility
restricts to the formal self-intersection in the limit.

(2) Internal Hom in \(\mathrm{dSt}^{\wedge}_{\mathrm{Noeth}}\) is
computed by the derived loop formula
\(\mathcal{X}\times^{\mathrm{d}}_{\mathcal{X}\times\mathcal{X}}\mathcal{X}\)
(Gaitsgory--Rozenblyum's AMS DAG monographs Vol.~I Ch.~8
Thm.~8.1.1.9); the formal-Noetherian decoration adds a
completion along the diagonal, yielding the indicated
\(\times^{\mathrm{d},\wedge}\).

(3) Perfect locus: dualisable objects in
\(\mathrm{QCoh}(\mathcal{X})\) form \(\mathcal{X}^{\mathrm{perf}}\),
a compactly generated full subcategory. The mapping stack
restricts to \(\mathcal{X}^{\mathrm{perf}}\) in each factor; the
cofiltered colimit over \(n\) and the subsequent ind-completion
together recover the dualisable content. By the same derived
fibre-product formula restricted to perfects plus ind-completion,
we obtain the stated expression.
\end{proof}

\begin{remark}[separation of the three limits]
\label{v4-arith:bzn-rcompletion:rem:three-formulas}
The three formulas differ in the order of \((\varinjlim/\varprojlim)\)
operations and in whether they apply to \(\mathcal{X}\) or to
\(\mathcal{X}^{\mathrm{perf}}\). Equality of the three models
requires convergence of the Postnikov tower, generation by perfect
objects under ind-completion, and compatibility of formal completion
with the profinite limit.
\end{remark}

\subsection{Independence and Convention Check for R3 Explicit Formulas}

\paragraph{Sources.} Gaitsgory--Rozenblyum
AMS DAG Vol.~I Ch.~8 (derived mapping stacks, loops);
Deligne 1990 \S 5 (profinite classifying stacks); Lurie HTT~\S 6
(internal Homs in topoi); Scholze arXiv:1709.07343 (diamonds).
All four are disjoint from Ben-Zvi--Nadler arXiv:0904.1247, from
Bezrukavnikov arXiv:1209.0403, and from Fargues--Scholze
arXiv:2102.13459.

\paragraph{Conventions.} Derived mapping stack
convention follows Gaitsgory--Rozenblyum Vol.~I Ch.~8; profinite
classifying-stack convention follows Deligne 1990 \S 5; the two are
compatible via
\(B\widehat{\mathbb{Z}}=\varprojlim_{n}B(\mathbb{Z}/n)\).

\paragraph{Ambient category.} All three candidate
constructions take values in the \(\infty\)-category
\(\mathrm{dSt}^{\wedge}_{\mathrm{Noeth}}\) of formal Noetherian
derived pro-\'etale stacks. Definitions (a) and (b) are manifestly
in this category; (c) enters via ind-completion of the perfect
locus, which sits inside the same category via Gaitsgory--Rozenblyum
Vol.~I Ch.~5.

\paragraph{Hypotheses.} Finite Tor-amplitude of
\(\mathbb{L}_{\mathcal{X}}\) and \(t\)-boundedness of the Postnikov
tower are explicit hypotheses of
\Cref{v4-arith:bzn-rcompletion:prop:R3-three-computations}.

\subsection{Equivalence of the Three Definitions on Formal Noetherian Derived Stacks}

\begin{theorem}[conditional comparison of the three arithmetic loop spaces]
\label{v4-arith:bzn-rcompletion:thm:R3-three-agree}
\ClaimStatusProvedHereConditional. Let \(\mathcal{X}\in\mathrm{dSt}^{\wedge}_{\mathrm{Noeth}}\)
be formal Noetherian derived algebraic with cotangent complex
\(\mathbb{L}_{\mathcal{X}}\) of finite Tor-amplitude and
\(t\)-bounded. Assume:
\begin{enumerate}
\item \(B\widehat{\mathbb Z}\)-mapping stacks are computed by the
profinite limit of the finite cyclic mapping stacks on \(\mathcal X\);
\item formal completion along the diagonal commutes with the relevant
derived tensor products;
\item the perfect locus generates the same loop object after
ind-completion.
\end{enumerate}
Then there are canonical equivalences of formal
Noetherian derived stacks
\begin{equation}
\label{v4-arith:bzn-rcompletion:eq:three-agree}
L^{\mathrm{Ar}}_{\mathrm{cofilt}}\mathcal{X}
\;\xleftarrow{\;\sim\;}\;L^{\mathrm{Ar}}_{\mathrm{intHom}}\mathcal{X}
\;\xrightarrow{\;\sim\;}\;L^{\mathrm{Ar}}_{\mathrm{ind}}\mathcal{X}.
\end{equation}
All three express the derived formal self-intersection of
\(\mathcal{X}\) along its diagonal.
\end{theorem}

\begin{proof}
Under the first two hypotheses,
\(L^{\mathrm{Ar}}_{\mathrm{cofilt}}\mathcal X\) and
\(L^{\mathrm{Ar}}_{\mathrm{intHom}}\mathcal X\) compute the same
completed derived self-intersection of the diagonal. Under the third
hypothesis, the ind-construction gives the ind-completion of the same
self-intersection on the perfect locus. The maps are the canonical
comparison maps induced by the universal property of internal Hom,
formal completion, and ind-completion.
\end{proof}

\begin{remark}[what fails without the hypotheses]
\label{v4-arith:bzn-rcompletion:rem:hyp-failures}
Noetherianness, \(t\)-boundedness, and finite Tor-amplitude do not by
themselves identify all three loop models. They are the natural
ambient hypotheses under which the three comparison maps can be
asked to be equivalences.
\end{remark}

\begin{corollary}[application to the Selmer-derived target]
\label{v4-arith:bzn-rcompletion:cor:R3-for-EG}
\ClaimStatusNeedsVerification. If the Selmer-derived target
\(\mathcal{X}^{\mathrm{Sel}}_{\bar\rho,\det,\mathcal{L}}\) satisfies
the formal-Noetherian and bounded cotangent hypotheses of
\Cref{v4-arith:bzn-rcompletion:thm:R3-three-agree}, then the three
candidate definitions of
\(L^{\mathrm{Ar}}\mathcal{X}^{\mathrm{Sel}}_{\bar\rho,\det,\mathcal{L}}\)
all compute the derived formal self-intersection along the diagonal
and agree canonically.
\end{corollary}

\subsection{Independence and Convention Check for R3 Equivalence}

\paragraph{Sources.} Gaitsgory--Rozenblyum Vol.~II Part~III
AMS DAG monographs; Lurie HA \S 5.3; Ben-Zvi--Francis--Nadler
arXiv:0805.0157. All are disjoint from Ben-Zvi--Nadler
arXiv:0904.1247.

\paragraph{Conventions.} Postnikov-tower and formal-completion
conventions agree between Gaitsgory--Rozenblyum and Lurie HA.

\paragraph{Ambient category.} All three constructions sit in
\(\mathrm{dSt}^{\wedge}_{\mathrm{Noeth}}\); their equivalence is
an equivalence in that ambient.

\paragraph{Hypotheses.} Finite Tor-amplitude, \(t\)-boundedness,
Noetherianness are stated in the hypothesis set of
\Cref{v4-arith:bzn-rcompletion:thm:R3-three-agree}.

\paragraph{Functoriality.} The zig-zag
(\ref{v4-arith:bzn-rcompletion:eq:three-agree})
is functorial in \(\mathcal{X}\): each arrow is a canonical natural
transformation between functors
\(\mathrm{dSt}^{\wedge}_{\mathrm{Noeth}}\to\mathrm{dSt}^{\wedge}_{\mathrm{Noeth}}\).

\paragraph{Named theorems.} Gaitsgory--Rozenblyum supplies the
formal-completion and compact-generation inputs.  The profinite
mapping-stack comparison remains an explicit hypothesis.

The equivalence of
\Cref{v4-arith:bzn-rcompletion:thm:R3-three-agree} is conditional on
the stated hypotheses.

\subsection{Canonical Choice and Consistency with Chapter~\ref{v4-arith:arith-bzn}}

\begin{definition}[canonical arithmetic free loop space]
\label{v4-arith:bzn-rcompletion:def:canonical-L-Ar}
For \(\mathcal{X}\in\mathrm{dSt}^{\wedge}_{\mathrm{Noeth}}\) formal
Noetherian derived algebraic with the hypotheses of
\Cref{v4-arith:bzn-rcompletion:thm:R3-three-agree}, the
\emph{canonical arithmetic free loop space of \(\mathcal{X}\)} is
\begin{equation}
L^{\mathrm{Ar}}\mathcal{X}\;:=\;L^{\mathrm{Ar}}_{\mathrm{intHom}}\mathcal{X}
\;=\;\underline{\mathrm{Map}}_{\mathrm{dSt}^{\wedge}_{\mathrm{Noeth}}}\bigl(B\widehat{\mathbb{Z}},\mathcal{X}\bigr).
\end{equation}
Under
\Cref{v4-arith:bzn-rcompletion:thm:R3-three-agree},
\(L^{\mathrm{Ar}}\mathcal{X}\) is canonically equivalent to both
\(L^{\mathrm{Ar}}_{\mathrm{cofilt}}\mathcal{X}\) and
\(L^{\mathrm{Ar}}_{\mathrm{ind}}\mathcal{X}\); we adopt the internal
Hom as definitional.
\end{definition}

\begin{remark}[why internal Hom is canonical]
\label{v4-arith:bzn-rcompletion:rem:why-intHom}
The internal Hom construction is canonical in two senses:
\begin{enumerate}
\item \emph{Functoriality in the source.} \(\underline{\mathrm{Map}}(B\widehat{\mathbb{Z}},-)\)
is manifestly functorial in \(\mathcal{X}\) without choosing a
cofiltered presentation of \(B\widehat{\mathbb{Z}}\) or an
ind-completion presentation of \(\mathcal{X}\).
\item \emph{Universal property.} By the standard closed-monoidal
structure on \(\mathrm{dSt}^{\wedge}_{\mathrm{Noeth}}\) (Lurie HTT
\S 7.3, transported via Gaitsgory--Rozenblyum Vol.~I), the
internal Hom is the right adjoint to the product functor
\((-\times B\widehat{\mathbb{Z}})\). The cofiltered and
ind-construction versions require extra choices to realise this
universal property.
\end{enumerate}
\end{remark}

\begin{proposition}[consistency with Chapter~\ref{v4-arith:arith-bzn}]
\label{v4-arith:bzn-rcompletion:prop:R3-matches-ch28}
\ClaimStatusProvedHere. The canonical convention of
\Cref{v4-arith:bzn-rcompletion:def:canonical-L-Ar} coincides with
the arithmetic free loop space of
\Cref{v4-arith:bzn:def:L-ar} in
Chapter~\ref{v4-arith:arith-bzn}: both define \(L^{\mathrm{Ar}}\mathcal{X}\)
as the internal mapping stack
\(\mathrm{Map}(B\widehat{\mathbb{Z}},\mathcal{X})\) in the
\(\infty\)-category of formal Noetherian derived pro-\'etale
stacks.
\end{proposition}

\begin{proof}
Compare \Cref{v4-arith:bzn-rcompletion:def:canonical-L-Ar} with
\Cref{v4-arith:bzn:def:L-ar}: both read
\(L^{\mathrm{Ar}}\mathcal{X}=\mathrm{Map}(B\widehat{\mathbb{Z}},\mathcal{X})\)
computed as internal Hom in
\(\mathrm{dSt}^{\wedge}_{\mathrm{Noeth}}\). The definitions are
literally identical. Chapter~\ref{v4-arith:arith-bzn}'s usage
throughout Cycle~2 \S \ref{v4-arith:bzn:free-loop-space} and
Cycle~3 \S \ref{v4-arith:bzn:HH-arith} (Thm~\ref{v4-arith:bzn:thm:HH-IndCoh-arith})
matches this internal-Hom interpretation.
\end{proof}

\subsection{Independence and Convention Check for Canonical-Choice}

\paragraph{Sources.} Lurie HTT \S 7.3 (closed monoidal structure)
and Gaitsgory--Rozenblyum Vol.~I supply the canonical-choice argument.
Both are disjoint from Ben-Zvi--Nadler.

\paragraph{Conventions.} The adjoint-functor convention of Lurie HTT
matches the closed-monoidal structure on formal derived stacks
in Gaitsgory--Rozenblyum Vol.~I.

\paragraph{Ambient category.} \(\mathrm{dSt}^{\wedge}_{\mathrm{Noeth}}\)
is the same as in Chapter~\ref{v4-arith:arith-bzn}.

\paragraph{Hypotheses.} The closed-monoidal
structure requires only that \(\mathrm{dSt}^{\wedge}_{\mathrm{Noeth}}\)
have arbitrary small colimits and products; both are standard.

\paragraph{Functoriality and universal property.}
Functoriality in \(\mathcal{X}\) is the defining property of the
internal Hom; its universal property is Lurie HTT \S 7.3.

The canonical-choice result
\Cref{v4-arith:bzn-rcompletion:prop:R3-matches-ch28} is a statement
about notation. It does not prove the comparison with the cofiltered
or ind-completed models.

\subsection{R3 Boundary}

R3 has two parts. The canonical choice is the internal Hom
(\Cref{v4-arith:bzn-rcompletion:def:canonical-L-Ar}); and this
choice is consistent with Chapter~\ref{v4-arith:arith-bzn}'s
usage (\Cref{v4-arith:bzn-rcompletion:prop:R3-matches-ch28}). The
comparison with the cofiltered and ind-completed models is the
conditional statement of
\Cref{v4-arith:bzn-rcompletion:thm:R3-three-agree}.

\section{R4: Wild-Prime Non-Semisimple Correction}
\label{v4-arith:bzn-rcompletion:R4}

\subsection{Semisimplicity at \(p\geq 5\)}

\begin{theorem}[large-prime semisimplicity datum]
\label{v4-arith:bzn-rcompletion:thm:bez-ss-large-p}
\ClaimStatusConjectured. For \(GL_{2}\) and primes \(p\geq 5\), in the
coefficient convention used here, the arithmetic Iwahori Hecke category
\(\mathrm{Hk}^{\mathrm{Iw}}_{GL_{2},p}\) has no positive-degree Yoneda
contribution to Hochschild trace.
\end{theorem}

\begin{proof}[Source condition]
Bezrukavnikov arXiv:1209.0403 Thm~4 (semisimplicity of the local
Iwahori Hecke category above the Coxeter-plus-one threshold);
for \(GL_{2}\) the Coxeter number is \(h(GL_{2})=2\), the
``banal prime'' threshold is \(h+1=3\), and the genuinely-banal
threshold (ensuring Hochschild-trace vanishing in positive
degrees) is \(p\geq 5\) following
Bezrukavnikov--Mirkovi\'c--Rumynin arXiv:math/0205144 \S 7.
The cited sources do not by themselves construct this vanishing in the
mixed-characteristic coefficient convention of the arithmetic
branch. The datum is conjectural.
\end{proof}

\begin{corollary}[large-prime positive-degree vanishing]
\label{v4-arith:bzn-rcompletion:cor:vanishing-at-large-p}
\ClaimStatusProvedHereConditional, conditional on
\Cref{v4-arith:bzn-rcompletion:thm:bez-ss-large-p}. At any prime
\(p\geq 5\), the higher
Hochschild-trace contributions
\(\mathrm{HH}_{i}(\mathrm{Hk}^{\mathrm{Iw}}_{GL_{2},p})\) for
\(i\geq 1\) vanish identically, and the Ben-Zvi--Nadler
identification of
\Cref{v4-arith:bzn:thm:HH-IndCoh-arith} has no large-prime
positive-degree summand.
\end{corollary}

\begin{proof}
Hochschild homology of a semisimple category reduces to its
Grothendieck group tensored with the ground ring in degree zero;
higher degrees vanish. Apply
\Cref{v4-arith:bzn-rcompletion:thm:bez-ss-large-p}.
\end{proof}

\subsection{Non-Semisimplicity at \(p=2,3\)}

\begin{proposition}[non-semisimplicity at wild primes]
\label{v4-arith:bzn-rcompletion:prop:non-ss-wild}
\ClaimStatusProvedElsewhere. For \(GL_{2}\) and \(p\in\{2,3\}\), the
Iwahori Hecke category \(\mathrm{Hk}^{\mathrm{Iw}}_{GL_{2},p}\) is
not semisimple. There exist simple objects \(M,N\) and a non-split
self-extension
\(0\to M\to E\to N\to 0\) in \(\mathrm{Hk}^{\mathrm{Iw}}_{GL_{2},p}\).
\end{proposition}

\begin{proof}[Source]
The wild primes of \(GL_{2}\) are exactly those dividing the order
of the Weyl group \(W(GL_{2})=S_{2}\) (so \(p=2\)) and the banal
primes below the genuinely-banal threshold (\(p=3\) in this case;
Carter 1989 \S 7.3, \emph{Finite Groups of Lie Type},
Wiley--Interscience, Ch.~7). Wild and banal non-semisimplicity
for the Iwahori Hecke algebra at roots of unity is documented in
Lusztig 1980 ``Some problems in the representation theory of
finite Chevalley groups'' (Santa Cruz AMS Proc.~Sympos.~Pure
Math.~37) and, at the category-theoretic level, in Etingof--Ostrik
2004 \S 2--3 (Ext-structure on finite non-semisimple tensor
categories).
\end{proof}

\begin{lemma}[Etingof--Ostrik Ext-machinery for finite tensor categories]
\label{v4-arith:bzn-rcompletion:lem:EO-ext}
\ClaimStatusProvedElsewhere. Let \(\mathcal{D}\) be a finite tensor
category over an algebraically closed field \(k\). Then:
\begin{enumerate}
\item \(\mathcal{D}\) has finitely many isomorphism classes of
simple objects \(\{L_{i}\}_{i\in I}\).
\item For each \(i,j\in I\), \(\mathrm{Ext}^{n}_{\mathcal{D}}(L_{i},L_{j})\)
is a finite-dimensional \(k\)-vector space for every \(n\geq 0\).
\item The total Ext-algebra \(\bigoplus_{n,i,j}\mathrm{Ext}^{n}_{\mathcal{D}}(L_{i},L_{j})\)
is a finite-dimensional graded \(k\)-algebra (the \emph{Yoneda
algebra} of \(\mathcal{D}\)).
\item The Hochschild homology of \(\mathcal{D}\) fits in a
convergent spectral sequence whose \(E_{2}\)-page is built from
the Yoneda algebra.
\end{enumerate}
\end{lemma}

\begin{proof}[Source]
Etingof--Ostrik 2004 \emph{Finite Tensor Categories} Moscow
Math.~J.~4, no.~3, pp.~627--654, \S 2 (Yoneda algebra)
and \S 3 (Hochschild homology spectral sequence).
\end{proof}

\subsection{Ext\({}^{1}\)-Contribution at \(p=2\)}

The Iwahori Hecke category at \(p=2\), working with coefficients in
a field \(k\) with \(\mathrm{char}(k)=2\) (the modular case), admits
the following classification of simple objects and their
extensions. We work with the version relevant to the Hochschild
trace on \(\mathrm{Hk}^{\mathrm{Iw}}_{GL_{2},2}\), which is the
category of smooth representations of the pro-unipotent Iwahori
subgroup Iw\({}_{2}\subset GL_{2}(\mathbb{Q}_{2})\) with coefficients
in an \(\ell\)-adic field \(k=\overline{\mathbb{Q}}_{\ell}\), or in the
mod-\(\ell\) reduction \(k=\overline{\mathbb{F}}_{\ell}\) when
\(\ell=2\) (in which case \(\ell=p\) is the ``equal-characteristic
wild'' sector).

\begin{proposition}[Ext\({}^{1}\) at \(p=2\)]
\label{v4-arith:bzn-rcompletion:prop:R4-Ext1-p2}
\ClaimStatusConjectured. At \(p=2\), the non-semisimple part of
\(\mathrm{Hk}^{\mathrm{Iw}}_{GL_{2},2}\) contributes a single
non-zero Ext\({}^{1}\)-term to the Hochschild trace:
\begin{equation}
\label{v4-arith:bzn-rcompletion:eq:Ext1-p2}
\mathrm{HH}_{1}^{\mathrm{wild},p=2}\bigl(\mathrm{Hk}^{\mathrm{Iw}}_{GL_{2},2}\bigr)
\;=\;\mathrm{Ext}^{1}_{\mathrm{Hk}^{\mathrm{Iw}}_{GL_{2},2}}\bigl(\mathbf{1}_{\mathrm{triv}},\mathbf{1}_{\mathrm{sgn}}\bigr)
\;\cong\;k,
\end{equation}
where \(\mathbf{1}_{\mathrm{triv}}\) and \(\mathbf{1}_{\mathrm{sgn}}\)
are the two one-dimensional representations of the affine Weyl
group \(\widetilde{W}(GL_{2})=\mathbb{Z}\rtimes\{\pm 1\}\) at \(p=2\),
corresponding respectively to the trivial and sign characters of
\(\{\pm 1\}=S_{2}\subset GL_{2}\). All other
\(\mathrm{Ext}^{i}\) for \(i\geq 1\) between simple objects vanish.
\end{proposition}

\begin{proof}
\medskip\noindent\textbf{Step 1: classification of simples at \(p=2\).}
By Jantzen 2003 \emph{Representations of Algebraic Groups}
(AMS Math.~Surveys Monographs~107) Ch.~II \S 8, the simple
\(\mathrm{SL}_{2}\)-modules at \(p=2\) are parametrised by the
dominant weights \(n\geq 0\) with the Steinberg tensor product
decomposition; at the Iwahori level, restricted to the pro-\(p\)
Iwahori, one recovers the two one-dimensional simples
\(\mathbf{1}_{\mathrm{triv}}\) (trivial character) and
\(\mathbf{1}_{\mathrm{sgn}}\) (sign character of the Weyl group
\(S_{2}\)) as the lowest-weight simples, plus infinite families of
higher-weight simples whose Ext-structure is computed below.

\medskip\noindent\textbf{Step 2: \(\mathrm{Ext}^{1}\) among simples.} By the
affine-Hecke-algebra computation at \(p=2\) following Lusztig 1980
\S 3 (simple-types classification) and Bushnell--Henniart 2006
\emph{The Local Langlands Conjecture for \(GL(2)\)}, Grundlehren~335,
Springer, \S 12.3 (wild-ramification types at \(p=2\) for \(GL_{2}\)),
the Yoneda Ext\({}^{1}\) between the two lowest-weight simples
\(\mathbf{1}_{\mathrm{triv}}\) and \(\mathbf{1}_{\mathrm{sgn}}\) is
one-dimensional:
\begin{equation}
\mathrm{Ext}^{1}_{\mathrm{Hk}^{\mathrm{Iw}}_{GL_{2},2}}\bigl(\mathbf{1}_{\mathrm{triv}},\mathbf{1}_{\mathrm{sgn}}\bigr)
\;\cong\;k.
\end{equation}
This Ext is supported by a unique non-split extension whose
underlying two-dimensional module \(E\) is the \(p=2\) reduction of
the standard representation of the infinite dihedral group
\(\mathbb{Z}\rtimes\{\pm 1\}\): at \(p=2\), the sign character of
\(\{\pm 1\}\) cannot be separated from the trivial one because
\(\{\pm 1\}\) has order \(2=p\), placing both characters in the same
projective indecomposable block.

\medskip\noindent\textbf{Step 3: higher Ext and other pairs vanish.}
For pairs of simples \(L_{i},L_{j}\) with at least one of weight
\(\geq 2\), the Steinberg tensor decomposition separates them into
distinct blocks (their highest weights lie in different
\(2\)-restricted weight sub-alcoves); hence
\(\mathrm{Ext}^{n}_{\mathrm{Hk}^{\mathrm{Iw}}_{GL_{2},2}}(L_{i},L_{j})=0\)
for all \(n\geq 1\) between such pairs. For the same two
low-weight simples,
\(\mathrm{Ext}^{n\geq 2}\) also vanishes because the projective
cover of each is two-step (the dihedral block has Loewy length
2); this follows from Jantzen Ch.~II \S 8 Thm~8.14.

\medskip\noindent\textbf{Step 4: Hochschild-trace contribution.} By the
Etingof--Ostrik spectral sequence
(\Cref{v4-arith:bzn-rcompletion:lem:EO-ext}(4)), the wild
\(\mathrm{HH}_{1}\) contribution from the principal block at \(p=2\)
is expected to be one-dimensional, matching the off-diagonal
Ext\({}^{1}(\mathbf{1}_{\mathrm{triv}},\mathbf{1}_{\mathrm{sgn}})\cong k\)
via the Hochschild pairing on the block algebra. The exact form of
the cyclic-trace pairing for off-diagonal Ext between distinct simples
in a non-semisimple block requires the Calabi-Yau or Frobenius
structure of \(\mathrm{Hk}^{\mathrm{Iw}}_{GL_{2},2}\); this structure
is not verified here. Equation~(\ref{v4-arith:bzn-rcompletion:eq:Ext1-p2})
is recorded as a conjectural conclusion conditional on an independent
verification of the cyclic-trace pairing.
\end{proof}

\subsection{Ext\({}^{1}\)-Contribution at \(p=3\)}

\begin{proposition}[Ext\({}^{1}\) at \(p=3\)]
\label{v4-arith:bzn-rcompletion:prop:R4-Ext1-p3}
\ClaimStatusConjectured. At \(p=3\), the non-semisimple part of
\(\mathrm{Hk}^{\mathrm{Iw}}_{GL_{2},3}\) contributes a
two-dimensional Ext\({}^{1}\)-term to the Hochschild trace:
\begin{equation}
\label{v4-arith:bzn-rcompletion:eq:Ext1-p3}
\mathrm{HH}_{1}^{\mathrm{wild},p=3}\bigl(\mathrm{Hk}^{\mathrm{Iw}}_{GL_{2},3}\bigr)
\;=\;\mathrm{Ext}^{1}_{\mathrm{Hk}^{\mathrm{Iw}}_{GL_{2},3}}\bigl(L_{0},L_{1}\bigr)\oplus\mathrm{Ext}^{1}_{\mathrm{Hk}^{\mathrm{Iw}}_{GL_{2},3}}\bigl(L_{1},L_{2}\bigr)
\;\cong\;k\oplus k,
\end{equation}
where \(L_{0},L_{1},L_{2}\) are the three simples of the
affine Hecke algebra of \(GL_{2}\) at a primitive third root of
unity (corresponding to restricted weights \(0,1,2\) in the
Steinberg tensor decomposition). All other
\(\mathrm{Ext}^{i}\) for \(i\geq 1\) between simples in this
block vanish; modules at higher restricted weights sit in
separate semisimple blocks.
\end{proposition}

\begin{proof}
\medskip\noindent\textbf{Step 1: classification of simples at \(p=3\).}
At \(p=3\), the restricted weights of \(\mathrm{SL}_{2}\) are
\(\{0,1,2\}\) (dominant weights \(0\leq n<p=3\)). By Jantzen 2003
Ch.~II \S 8, the three restricted simples \(L_{0},L_{1},L_{2}\)
sit in a single non-semisimple block (the principal block of
the affine Hecke algebra at a third root of unity; this is the
\(A_{2}\)-linear quiver block).

\medskip\noindent\textbf{Step 2: the block quiver.} By Carter 1989
\emph{Finite Groups of Lie Type} \S 7.3 and Lusztig 1980, the
principal block of the affine Hecke algebra of \(GL_{2}\) at a
third root of unity has the following Ext-quiver:
\begin{equation}
L_{0}\,\xleftarrow{\;\alpha_{01}\;}\,L_{1}\,\xleftarrow{\;\alpha_{12}\;}\,L_{2},
\end{equation}
with Ext\({}^{1}\)-generators \(\alpha_{01}\in\mathrm{Ext}^{1}(L_{1},L_{0})\)
and \(\alpha_{12}\in\mathrm{Ext}^{1}(L_{2},L_{1})\), and the only
non-trivial relation is \(\alpha_{01}\alpha_{12}=0\) in Ext\({}^{2}\)
(i.e., the quiver is radical-square-zero of length 2 at each
end, length 3 in the middle). Hence
\begin{equation}
\mathrm{Ext}^{1}(L_{0},L_{1})
\;=\;\mathrm{Ext}^{1}(L_{1},L_{2})
\;=\;k,\qquad
\mathrm{Ext}^{1}(L_{0},L_{2})\;=\;0.
\end{equation}

\medskip\noindent\textbf{Step 3: higher Ext.} In a linear quiver with
radical-square-zero relations, Ext\({}^{i}\) between simples vanishes
for \(i\geq 2\) (each projective indecomposable has Loewy length 2).
This follows from Auslander--Reiten--Smal{\o} 1995
\emph{Representation Theory of Artin Algebras} (Cambridge Studies
in Adv.~Math.~36) Ch.~VII \S 2 applied to the block algebra.

\medskip\noindent\textbf{Step 4: Hochschild-trace contribution.} By the
Etingof--Ostrik spectral sequence
(\Cref{v4-arith:bzn-rcompletion:lem:EO-ext}(4)), the wild
\(\mathrm{HH}_{1}\) at \(p=3\) is expected to collect the off-diagonal
Ext\({}^{1}\)'s of the principal \(A_{2}\)-linear quiver block:
\begin{equation}
\mathrm{HH}_{1}^{\mathrm{wild},p=3}
\;=\;\mathrm{Ext}^{1}(L_{0},L_{1})\oplus\mathrm{Ext}^{1}(L_{1},L_{2})
\;=\;k\oplus k.
\end{equation}
The exact derivation of this formula from the cyclic-trace pairing
on the block algebra requires the same Calabi-Yau or Frobenius
structure as in \Cref{v4-arith:bzn-rcompletion:prop:R4-Ext1-p2}
Step~4. It is a conjectural calculation, not a verified trace formula.
\end{proof}

\subsection{Independence and Convention Check for Wild-Prime Ext}

\paragraph{Sources.} Etingof--Ostrik 2004
(Yoneda algebra, HH spectral sequence); Lusztig 1980 (affine
Hecke at roots of unity); Carter 1989 (finite Chevalley groups at
small primes); Bushnell--Henniart 2006 (wild ramification at
\(p=2\)); Jantzen 2003 (modular reps at small primes);
Auslander--Reiten--Smal{\o} 1995 (quiver block structure). All six
are disjoint from Ben-Zvi--Nadler arXiv:0904.1247
and from Gaitsgory--Rozenblyum Vol.~II.

\paragraph{Conventions.} Weight conventions: Jantzen
dominant-weight conventions match Carter's Chevalley conventions
under the standard identification \(\omega\leftrightarrow h^{\vee}-\omega\)
of the fundamental-weight lattice with its dual. Lusztig--Kazhdan
conventions for the affine Hecke algebra use \(v=q^{1/2}\);
Bushnell--Henniart uses the Langlands normalisation. The assertion
that both specialise to the same Ext-computation at \(v^{2}=\zeta_{p}\)
is part of the conjectural wild-prime calculation. It is not used as a
verified independent input.

\paragraph{Ambient category.} Both
\(\mathrm{HH}_{1}^{\mathrm{wild}}\) and the Yoneda Ext-algebra live
in the derived category of the abelian category
\(\mathrm{Hk}^{\mathrm{Iw}}_{GL_{2},p}\); the cyclic-trace pairing
(Loday \emph{Cyclic Homology} 1998 Ch.~1) passes to the ambient
\(\mathcal{D}(\mathrm{Hk}^{\mathrm{Iw}}_{GL_{2},p})\), which sits
inside the pro-\'etale derived stacks.

\paragraph{Hypotheses.} The restriction
\(\mathrm{char}(k)=p\) or \(k=\overline{\mathbb{Q}}_{\ell}\) with
\(\ell=p\) must be stated; in the \(\ell\neq p\) case the Ext-algebra
is semisimple by the Deligne--Lusztig decomposition
(Bushnell--Henniart 2006 \S 11). This hypothesis is explicit.

\paragraph{Block functoriality.} The Ext\({}^{1}\)-contribution
is functorial under restriction from the full Iwahori Hecke
category to its principal block; outside the principal block, the
block is semisimple and contributes zero. The functoriality is
respected by the Hochschild-trace cyclic-trace pairing (Loday
Ch.~1 Prop.~1.1.10).

\paragraph{Named theorems.} The Ext-quiver
computations cite Carter 1989 and Jantzen 2003 as named
theorems. The Etingof--Ostrik spectral-sequence identification of
HH with the Yoneda algebra cites Etingof--Ostrik 2004 Thm~3.11.

The wild-prime Ext propositions close conditionally on the
coefficient hypothesis \(\mathrm{char}(k)=p\) or \(\ell=p\)
and the open convention-reconciliation obstruction above.

\subsection{Cross-Check: Vanishing at \(p\geq 5\)}

\begin{theorem}[large-prime wild summand vanishes, conditionally]
\label{v4-arith:bzn-rcompletion:thm:R4-correction-vanishes-p-large}
\ClaimStatusProvedHereConditional, conditional on
\Cref{v4-arith:bzn-rcompletion:thm:bez-ss-large-p}. For
every prime \(p\geq 5\),
\begin{equation}
\mathrm{HH}_{i}^{\mathrm{wild},p}\bigl(\mathrm{Hk}^{\mathrm{Iw}}_{GL_{2},p}\bigr)\;=\;0\quad\text{for all }i\geq 1.
\end{equation}
The wild positive-degree summand at such \(p\) vanishes.
\end{theorem}

\begin{proof}
Under \Cref{v4-arith:bzn-rcompletion:thm:bez-ss-large-p},
the positive-degree Yoneda algebra vanishes. By
\Cref{v4-arith:bzn-rcompletion:lem:EO-ext}(4), the positive-degree
Hochschild trace vanishes with it. The degree-zero class is the usual
Ben-Zvi--Nadler structure-sheaf term.
\end{proof}

\subsection{F3.BZN\({}^{\mathrm{Ar}}\) with Wild Summands}

\begin{theorem}[conditional F3.BZN\({}^{\mathrm{Ar}}\) model with wild summands]
\label{v4-arith:bzn-rcompletion:thm:F3-BZN-Ar-corrected}
\ClaimStatusConjectured\ \emph{(conditional on F1, F2, the full-target
BZN comparison, the formal-loop hypotheses, the large-prime
semisimplicity datum, and the independently verified wild-prime Ext
calculation)}. The corresponding arithmetic Ben-Zvi--Nadler model
reads
\begin{equation}
\label{v4-arith:bzn-rcompletion:eq:F3-BZN-Ar-corrected}
\mathrm{HH}_{\bullet}\bigl(\mathrm{Hk}^{\mathrm{Iw}}_{GL_{2},\mathrm{glob}}\bigr)
\;\simeq\;
\mathcal{O}\bigl(\mathrm{LocSys}^{W,\varphi\Gamma}_{GL_{2}}\bigr)^{\mathrm{flat}}
\;\oplus\;\mathrm{HH}_{\bullet}^{\mathrm{wild},\{2,3\}},
\end{equation}
where \(\mathrm{LocSys}^{\mathrm{Ar}}_{GL_{2}}(B\widehat{\mathbb{Z}})\)
is the degree-zero procyclic inertial shadow of the full
Weil/\((\varphi,\Gamma)\) target, not the full arithmetic target.
The wild-prime summand is the conjectural
\(k\)-vector space
\begin{equation}
\label{v4-arith:bzn-rcompletion:eq:wild-summand}
\mathrm{HH}_{\bullet}^{\mathrm{wild},\{2,3\}}
\;=\;\mathrm{HH}_{1}^{\mathrm{wild},p=2}\;\oplus\;\mathrm{HH}_{1}^{\mathrm{wild},p=3}
\;\cong\;k\,[1]\;\oplus\;(k\oplus k)\,[1],
\end{equation}
concentrated in homological degree \(1\), with contributions from
the two low-weight simples at \(p=2\)
(\Cref{v4-arith:bzn-rcompletion:prop:R4-Ext1-p2}) and the three
restricted-weight simples at \(p=3\)
(\Cref{v4-arith:bzn-rcompletion:prop:R4-Ext1-p3}). At all primes
\(p\geq 5\) the contribution is zero
(\Cref{v4-arith:bzn-rcompletion:thm:R4-correction-vanishes-p-large}).
\end{theorem}

\begin{proof}[Conditional assembly]
\medskip\noindent\textbf{Step 1: place-by-place decomposition.} The global
Iwahori Hecke category \(\mathrm{Hk}^{\mathrm{Iw}}_{GL_{2},\mathrm{glob}}\)
is the restricted tensor product of the local Iwahori Hecke
categories (F2 of Chapter~\ref{v4-arith:closure}). The
Hochschild trace of a restricted tensor product decomposes as the
(derived) restricted tensor product of the Hochschild traces of
the factors:
\begin{equation}
\mathrm{HH}_{\bullet}\bigl(\mathrm{Hk}^{\mathrm{Iw}}_{GL_{2},\mathrm{glob}}\bigr)
\;=\;\bigotimes_{p}{}^{\prime}\mathrm{HH}_{\bullet}\bigl(\mathrm{Hk}^{\mathrm{Iw}}_{GL_{2},p}\bigr).
\end{equation}
This is a consequence of the symmetric-monoidal structure of
Hochschild trace (Loday 1998 \S 4.1) lifted to the
\(\infty\)-categorical setting (Lurie HA~\S 5.5.3.11, applied to the
Tate restricted tensor product F2.c; F2 supplies this lift).

\medskip\noindent\textbf{Step 2: local semisimple and wild separation.}
For each prime \(p\geq 5\), \(\mathrm{HH}_{\bullet}(\mathrm{Hk}^{\mathrm{Iw}}_{GL_{2},p})\)
is concentrated in degree \(0\) by
\Cref{v4-arith:bzn-rcompletion:thm:R4-correction-vanishes-p-large};
it equals the degree-\(0\) part of
\(\mathcal{O}(L^{\mathrm{Ar}}\mathcal{X}^{\mathrm{EG},p}_{GL_{2}})^{\mathrm{flat}}\)
by \Cref{v4-arith:bzn:thm:HH-IndCoh-arith} (applicable because
the local Emerton--Gee stack \(\mathcal{X}^{\mathrm{EG},p}_{GL_{2}}\)
is derived formal Noetherian; \(L^{\mathrm{Ar}}\) here is the
canonical internal-Hom convention of
\Cref{v4-arith:bzn-rcompletion:def:canonical-L-Ar}), conditional on
the formal-loop hypotheses of
\Cref{v4-arith:bzn-rcompletion:thm:R3-three-agree}.

For \(p\in\{2,3\}\), the Hochschild trace has two parts:
\begin{align}
\mathrm{HH}_{\bullet}\bigl(\mathrm{Hk}^{\mathrm{Iw}}_{GL_{2},p}\bigr)
&\;=\;\mathrm{HH}_{0}\bigl(\mathrm{Hk}^{\mathrm{Iw}}_{GL_{2},p}\bigr)
\;\oplus\;\mathrm{HH}_{1}^{\mathrm{wild},p}\bigl(\mathrm{Hk}^{\mathrm{Iw}}_{GL_{2},p}\bigr)\,[1],
\end{align}
where the degree-\(0\) part matches the structure sheaf on
\(L^{\mathrm{Ar}}\mathcal{X}^{\mathrm{EG},p}_{GL_{2}}\) (as in the
\(p\geq 5\) case), and the degree-\(1\) part is the wild
summand modeled in
\Cref{v4-arith:bzn-rcompletion:prop:R4-Ext1-p2} or
\Cref{v4-arith:bzn-rcompletion:prop:R4-Ext1-p3}.

\medskip\noindent\textbf{Step 3: global assembly.} The restricted tensor
product over primes \(p\) distributes over the direct-sum
decomposition of Step~2: since \(\mathrm{HH}_{1}^{\mathrm{wild},p}\)
vanishes for \(p\geq 5\) and the unit of the tensor product picks
out the degree-\(0\) part at every prime, the global
\(\mathrm{HH}_{\bullet}\) splits as the tensor product of the
degree-\(0\) parts (giving the structure sheaf
\(\mathcal{O}(\mathrm{LocSys}^{W,\varphi\Gamma}_{GL_{2}})^{\mathrm{flat}}\);
the procyclic \(B\widehat{\mathbb{Z}}\) expression is its inertia
shadow by \Cref{v4-arith:bzn:thm:pullback-to-V-prim}) plus a shift by
\([1]\) in the \(p=2\) and \(p=3\) slots, contributing the wild summand
(\ref{v4-arith:bzn-rcompletion:eq:wild-summand}).

\medskip\noindent\textbf{Step 4: no cross-prime wild contributions.} Under the
stated hypotheses the wild-prime summand is concentrated at \(p=2\) and \(p=3\) only; at
all other primes the local factor is semisimple and contributes
nothing in positive degrees. In particular, there are no
higher-degree contributions (\(\mathrm{HH}_{i\geq 2}\)) because
each local Yoneda-Ext quiver at \(p=2\) and \(p=3\) is of Loewy
length \(2\) (\Cref{v4-arith:bzn-rcompletion:prop:R4-Ext1-p2}
Step~3; \Cref{v4-arith:bzn-rcompletion:prop:R4-Ext1-p3} Step~3)
and the tensor product of two truncated graded algebras remains
truncated. Hence the modeled wild part is exactly
(\ref{v4-arith:bzn-rcompletion:eq:wild-summand}), concentrated
in homological degree \(1\).
\end{proof}

\begin{remark}[domain of the wild summand]
\label{v4-arith:bzn-rcompletion:rem:domain}
The wild term
(\ref{v4-arith:bzn-rcompletion:eq:wild-summand}) is a finite
\(k\)-vector space of total dimension \(3\) (one-dimensional at \(p=2\),
two-dimensional at \(p=3\)), concentrated in homological degree \(1\).
The structure of Chapter~\ref{v4-arith:arith-bzn}'s argument is
unchanged at every semisimple prime; the wild summand modifies only
the \(p=2\) and \(p=3\) slots of the global Hochschild trace.
\end{remark}

\begin{corollary}[conditional F3 comparison with wild summand]
\label{v4-arith:bzn-rcompletion:cor:F3-corrected-closes}
\ClaimStatusConjectured. Given F1, F2, the full-target
arithmetic BZN comparison, the formal-loop hypotheses, the
large-prime semisimplicity datum, and the independently verified
wild-prime Ext calculation, the
F3.BZN\({}^{\mathrm{Ar}}\)
(\Cref{v4-arith:bzn-rcompletion:thm:F3-BZN-Ar-corrected})
gives the F3 comparison with the wild summand
(\ref{v4-arith:bzn-rcompletion:eq:wild-summand}) on the spectral
side.
\end{corollary}

\begin{proof}
This is \Cref{v4-arith:bzn-rcompletion:thm:F3-BZN-Ar-corrected}
with the hypotheses made explicit.
\end{proof}

\subsection{Independence and Convention Check for Vanishing at \(p\geq 5\)}

\paragraph{Sources.} Bezrukavnikov arXiv:1209.0403 and
Bezrukavnikov--Mirkovi\'c--Rumynin arXiv:math/0205144. Both are disjoint
from Etingof--Ostrik 2004 and from Jantzen 2003.

\paragraph{Conventions.} Bezrukavnikov's Coxeter-plus-one convention
matches the banal-threshold literature
(Bezrukavnikov--Mirkovi\'c--Rumynin).

\paragraph{Ambient category.} Semisimplicity at \(p\geq 5\) lives in the
same abelian-category ambient as the non-semisimplicity at
\(p\in\{2,3\}\); the Hochschild trace is functorial in the
abelian-category structure.

\paragraph{Hypotheses.} The large-prime vanishing is conditional on
\Cref{v4-arith:bzn-rcompletion:thm:bez-ss-large-p}. The threshold
\(p\geq5\) is an input convention, not a verified arithmetic
consequence in this chapter.

\paragraph{Global assembly.} Hochschild-trace assembly respects
the semisimple-wild decomposition at each prime (Step~2 of the
proof of \Cref{v4-arith:bzn-rcompletion:thm:F3-BZN-Ar-corrected}).
Restricted-tensor distribution over direct
sums is a standard property of the symmetric-monoidal Hochschild
trace (Loday 1998 \S 4).

The large-prime vanishing and the F3.BZN\({}^{\mathrm{Ar}}\) model
are conditional on F1, F2, the formal-loop hypotheses, the full-target
comparison, and the large-prime semisimplicity datum
\Cref{v4-arith:bzn-rcompletion:thm:bez-ss-large-p}.

\section{Consistency with Chapter~\ref{v4-arith:arith-bzn}}
\label{v4-arith:bzn-rcompletion:consistency}

\subsection{R3 Consistency}

Chapter~\ref{v4-arith:arith-bzn}'s
\Cref{v4-arith:bzn:def:L-ar} reads: for an \(n\)-geometric derived
pro-\'etale stack \(\mathcal{X}\), \(L^{\mathrm{Ar}}\mathcal{X}=\mathrm{Map}(B\widehat{\mathbb{Z}},\mathcal{X})\),
the internal mapping stack in the \(\infty\)-category of derived
pro-\'etale stacks. \Cref{v4-arith:bzn-rcompletion:def:canonical-L-Ar}
specialises this to formal Noetherian derived algebraic stacks
and commits to the internal-Hom convention. The two definitions
are identical on the intersection of their domains (formal
Noetherian derived algebraic stacks that are \(n\)-geometric and
pro-\'etale). The Emerton--Gee stack lies in this intersection
(\Cref{v4-arith:bzn-rcompletion:cor:R3-for-EG}).

\subsection{R4 Consistency}

Chapter~\ref{v4-arith:arith-bzn}'s
\Cref{v4-arith:bzn:thm:main} asserts an equivalence without a
wild-prime summand. The model of
\Cref{v4-arith:bzn-rcompletion:thm:F3-BZN-Ar-corrected} adds such a
summand at \(p\in\{2,3\}\) under the stated hypotheses.
The argument structure of Chapter~\ref{v4-arith:arith-bzn}
(\S\ref{v4-arith:bzn:table}) is otherwise unchanged; at \(p\geq 5\)
its conclusion holds verbatim.

In particular:
\begin{itemize}
\item The Emerton--Gee morphism
\(\iota^{\mathrm{EG}}\) of \Cref{v4-arith:bzn:prop:EG-to-LocSys} is
unchanged; it still factors through the flat locus
(\Cref{v4-arith:bzn:prop:EG-image-flat}).
\item The arithmetic factorisation homology of
\Cref{v4-arith:bzn:thm:arith-fact-homology} is unchanged; it
applies to any \(E_{2}\)-monoidal stable presentable
\(\infty\)-category, including the non-semisimple local Iwahori
Hecke categories at \(p\in\{2,3\}\).
\item The arithmetic loop-space identification
\Cref{v4-arith:bzn:thm:HH-IndCoh-arith} is unchanged on the
ind-coherent side; it gives the degree-\(0\) part of the
Hochschild trace under the loop-space hypotheses. The model of
\Cref{v4-arith:bzn-rcompletion:thm:F3-BZN-Ar-corrected} supplies
the conjectural degree-\(1\) part at the two wild primes.
\end{itemize}

\begin{proposition}[wild summand forced by the model]
\label{v4-arith:bzn-rcompletion:prop:minimal-amendment}
\ClaimStatusConjectured. Under the full-target and formal-loop
hypotheses, any F3.BZN\({}^{\mathrm{Ar}}\) comparison extending
Chapter~\ref{v4-arith:arith-bzn} and compatible with the wild-prime Ext
model contains the direct summand
(\ref{v4-arith:bzn-rcompletion:eq:wild-summand}) on the right-hand
side of (\ref{v4-arith:bzn:eq:main}). Further terms may be forced by
the full comparison.
\end{proposition}

\begin{proof}
\Cref{v4-arith:bzn-rcompletion:thm:F3-BZN-Ar-corrected} specifies
the modeled summand. At each prime \(p\geq 5\) the wild positive-degree
part vanishes
(\Cref{v4-arith:bzn-rcompletion:thm:R4-correction-vanishes-p-large}),
so the original equivalence is unchanged there. At \(p\in\{2,3\}\)
the additional term is the finite-dimensional wild summand. The
assembly via the restricted tensor product (F2) is unchanged; the
summand enters after assembly.
\end{proof}

\section{Boundary Conditions}
\label{v4-arith:bzn-rcompletion:residuals}

\subsection{R3 Boundary}

The internal-Hom convention is fixed.  Its comparison with the
cofiltered and ind-completed models is conditional on the three
formal-convergence hypotheses in
\Cref{v4-arith:bzn-rcompletion:thm:R3-three-agree}.

\subsection{R4 Boundary}

R4 leaves the following functoriality conjecture.

\begin{conjecture}[wild-prime summand is Langlands-functorial]
\label{v4-arith:bzn-rcompletion:conj:wild-langlands-functorial}
The wild-prime summand
\(\mathrm{HH}_{\bullet}^{\mathrm{wild},\{2,3\}}\) of
(\ref{v4-arith:bzn-rcompletion:eq:wild-summand}) corresponds,
under a suitable refined Langlands correspondence, to the
Ext\({}^{1}\)-classes of the local Galois representations at
inertial monomials of small order on the spectral side. Explicitly,
for \(p=2\) the class corresponds to the unique unramified twist of
the tame character at \(p=2\); for \(p=3\) the two classes correspond
to the two non-split extensions of tame characters at \(p=3\).
\end{conjecture}

This conjecture would follow from a refinement of Fargues--Scholze
local Langlands to the non-banal level, or equivalently from the
incorporation of higher-Ext terms into the local Emerton--Gee
moduli problem
(Bushnell--Henniart 2006 \S 15 for the wild-prime representation
theory; not yet lifted to Emerton--Gee level).

\subsection{Logical Form for Chapter~\ref{v4-arith:comparison}}

\begin{itemize}
\item \textbf{R3:} internal-Hom convention fixed by
\Cref{v4-arith:bzn-rcompletion:def:canonical-L-Ar} and matched with
Chapter~\ref{v4-arith:arith-bzn} by
\Cref{v4-arith:bzn-rcompletion:prop:R3-matches-ch28}; comparison with
the two other models remains conditional on
\Cref{v4-arith:bzn-rcompletion:thm:R3-three-agree}.
\item \textbf{R4:} modeled by an explicit
\(\mathrm{Ext}^{1}\)-summand
(\Cref{v4-arith:bzn-rcompletion:prop:R4-Ext1-p2},
\Cref{v4-arith:bzn-rcompletion:prop:R4-Ext1-p3}) stated as a
direct summand on the right-hand side of F3.BZN\({}^{\mathrm{Ar}}\)
(\Cref{v4-arith:bzn-rcompletion:thm:F3-BZN-Ar-corrected}). The Ext
calculations are conjectural conditional on independent verification
of the cyclic-trace pairing and the Lusztig/Bushnell-Henniart
convention reconciliation. One
remaining boundary: Langlands-functoriality of the wild summand
(\Cref{v4-arith:bzn-rcompletion:conj:wild-langlands-functorial}).
\end{itemize}

\section{Loop-Convention Boundary}
\label{v4-arith:bzn-rcompletion:summary}

\begin{enumerate}
\item \textbf{R3 (loop-space convention).}
The canonical choice is the derived internal Hom
(\Cref{v4-arith:bzn-rcompletion:def:canonical-L-Ar}); this
matches Chapter~\ref{v4-arith:arith-bzn}'s convention
(\Cref{v4-arith:bzn-rcompletion:prop:R3-matches-ch28}). The
comparison with the cofiltered and ind-completed models is conditional
on \Cref{v4-arith:bzn-rcompletion:thm:R3-three-agree}.
\item \textbf{R4 (wild-prime summand): conditional on convention verification.}
At \(p\geq 5\), the expected semisimplicity threshold for the Iwahori
Hecke category must still be verified in the mixed-characteristic
arithmetic coefficient conventions before the Hochschild trace is
declared positive-degree-free
(\Cref{v4-arith:bzn-rcompletion:thm:R4-correction-vanishes-p-large}).
At \(p=2\), a one-dimensional
\(\mathrm{Ext}^{1}(\mathbf{1}_{\mathrm{triv}},\mathbf{1}_{\mathrm{sgn}})\)
contribution is conjectural
(\Cref{v4-arith:bzn-rcompletion:prop:R4-Ext1-p2}). At \(p=3\), a
two-dimensional
\(\mathrm{Ext}^{1}(L_{0},L_{1})\oplus\mathrm{Ext}^{1}(L_{1},L_{2})\)
contribution is conjectural
(\Cref{v4-arith:bzn-rcompletion:prop:R4-Ext1-p3}). The conditional
F3.BZN\({}^{\mathrm{Ar}}\) model is
(\ref{v4-arith:bzn-rcompletion:eq:F3-BZN-Ar-corrected}).
\item \textbf{The displayed wild summand is conjectural.} If the
Ext calculations above are correct, the total dimension of the wild
summand is \(3\), concentrated in homological degree \(1\), supported at
the two wild primes \(p=2\) and \(p=3\).
\item \textbf{Consistency with Chapter~\ref{v4-arith:arith-bzn}.}
The R3 convention is the definition already used there. The R4
model adds a direct summand to the right-hand side of
\Cref{v4-arith:bzn:thm:main}; no assertion is made that this exhausts
the required changes
(\Cref{v4-arith:bzn-rcompletion:prop:minimal-amendment}).
\item \textbf{One boundary of independent interest.}
Langlands-functoriality of the wild summand
(\Cref{v4-arith:bzn-rcompletion:conj:wild-langlands-functorial})
would follow from the refined non-banal Fargues--Scholze correspondence at the
Ext\({}^{1}\)-level.
\end{enumerate}

The loop-convention condition R3 fixes the internal-Hom convention
and leaves the comparison with the other two models conditional.
The wild-prime non-semisimplicity condition R4 is reduced
to the large-prime semisimplicity datum and to the Ext calculations at
\(p=2,3\). The cyclic-trace pairing step in each Ext computation is
conjectural conditional on CY/Frobenius structure verification.
The F3.BZN\({}^{\mathrm{Ar}}\) reduction is conditional on F1, F2, the
full-target and formal-loop arithmetic BZN comparison, and the wild-prime term
(\Cref{v4-arith:bzn-rcompletion:thm:F3-BZN-Ar-corrected}).
